\section{Technical bounds}\label{app:technical}
\subsection{Controlled finite-time Taylor remainders}
For a complexity written in its finite cyclic space as
$C(t)=\bra v e^{itL}N e^{-itL}\ket v$, with $N\ket v=0$, repeated
commutators give
\begin{equation}\label{eq:taylor-norm}
 |C'''(t)|\le(2\norm L)^3\norm N,\qquad
 |C(t)-kt^2|\le\frac{(2\norm L)^3\norm N}{6}|t|^3,
\end{equation}
where $k$ is the quadratic coefficient. For Eq.~\eqref{eq:right-state},
both Liouvillian norms are at most two and the number-operator norms are
two and four. Adding the two remainder estimates yields $64|t|^3$,
proving Eq.~\eqref{eq:right-interval}.

For $H_\delta$ at $\delta=1/20$, the maximal absolute row sum bounds
$\norm{H_\delta}\le11/10$. Each operator chain has at most seven
distinct energy-gap atoms, so each number-operator norm is at most six.
The two Liouvillian norms are at most $11/5$. Hence
\begin{equation}\label{eq:perturbed-remainder}
 \CM(t)-\KI(t)
 \ge\frac{53}{3380}t^2-\frac{21296}{125}|t|^3.
\end{equation}
At $t=1/20000$ this is exactly the positive rational lower bound in
Eq.~\eqref{eq:robust-gap}. These are conservative finite-time bounds,
not numerical estimates of the exact differences.

\subsection{Period integral and mean estimates}
For $a,b>0$, tangent substitution after splitting the interval at $\pi$
gives
\begin{equation}\label{eq:integral}
 \frac1{2\pi}\int_0^{2\pi}
 \frac{dt}{a\cos^2(t/2)+b}
 =\frac1{\sqrt{b(a+b)}}.
\end{equation}
For the main-ratio family, $D_0\le1$, $1-\eps\ge3/4$, and
$\sqrt{4(1-\eps)^2+\eps^2}\ge2(1-\eps)$ bound the mean in
Eq.~\eqref{eq:main-mean} by $68/81$. On the peak window, its rational
denominator is at most $2\eps^2$ while its numerator is at least
$\eps$; together with $D_0\ge1/2$ this proves
Eq.~\eqref{eq:main-window}.

For the reciprocal family, the positive coefficients in
Eq.~\eqref{eq:reciprocal-traces} show that $2\le N,T_4\le3$ on
$0<\eps\le1/16$. The bracket in Eq.~\eqref{eq:reciprocal-mean} is at
most $9/2$, so its prefactor bounds the mean by $27\eps$.
On the window the denominator $4c+\eps^2$ is at most $2\eps^2$
and the numerator $16c+\eps$ is at least $\eps$. This gives
$g_\eps\ge2N/T_4\ge4/3$ and hence the second-moment estimate.

\subsection{The fixed-purity endpoint bound}
In Eq.~\eqref{eq:purity-ratio} the fractional-linear factor decreases
because $\eps/5>18\eps^2/25$ for $0<\eps<5/18$. Its value at
$x=0$ is one. At $x=1$ it is
\begin{equation}\label{eq:purity-endpoint}
 \frac{1-\eps/10}{1-9\eps^2/25}\ge1-\eps/10.
\end{equation}
This gives the lower bound in Eq.~\eqref{eq:purity-bounds}.
